\chapter{Relational Inference Histories and Fisher Information Clocks}
\label{ch:relational-inference}

The contextual base $\mathcal C$ is fixed data in the present theory.  It
organizes the domains on which agents are represented, but it carries no
distinguished temporal parameter.  Inference therefore has to be described
as change in a statistical fiber or as change of a section, rather than as
motion of the base.  A curve in $\mathcal C$ orders contexts and can be used
to probe a section; it is not, by that fact, an agent history.
\status{HYPOTHESIS}

This chapter supplies two structures that the preceding pullback construction
deliberately left open.  A declared VFE gradient ray selects an oriented
unparameterized inference orbit, and the Fisher line element assigns length
to that already selected orbit.  The resulting information clock is
agent-relative and invariant under orientation-preserving
reparameterization.  It is neither a physical time coordinate nor a scalar
that automatically exists on the whole configuration space.
\status{DEFINITION}

\section{Fixed-base kinematics}
\label{sec:hist-fixed-base-kinematics}

For agent $i$, restrict the belief and model bundles to its contextual domain
$\mathcal C_i$ and form
\begin{equation}
\mathcal E_i
=\mathcal E_b|_{\mathcal C_i}
\times_{\mathcal C_i}
\mathcal E_m|_{\mathcal C_i},
\qquad
\varpi_i:\mathcal E_i\longrightarrow\mathcal C_i.
\label{eq:hist-agent-product-bundle}
\end{equation}
The pair $\boldsymbol\omega_i=(\omega_b,\omega_m)$ supplies a product
horizontal distribution and a vertical projection
$\operatorname{ver}^{\boldsymbol\omega_i}$.  Its vertical Fisher form is not
fixed until positive channel weights or a joint statistical metric are
declared.  The common bundle projection supplies neither choice.
\status{DEFINITION}

\definitionheading{Vertical, horizontal, and mixed curves}{def:hist-curve-types}
Let $J$ be a connected interval and let
$\Gamma:J\to\mathcal E_i$ be $C^1$, with
$\gamma=\varpi_i\circ\Gamma$.  Relative to
$\boldsymbol\omega_i$, its velocity has the unique splitting
\begin{equation}
\dot\Gamma
=H^{\boldsymbol\omega_i}_{\Gamma}\dot\gamma
+v_{\boldsymbol\omega_i}(\Gamma),
\qquad
v_{\boldsymbol\omega_i}(\Gamma)
=\operatorname{ver}^{\boldsymbol\omega_i}(\dot\Gamma).
\label{eq:hist-total-curve-splitting}
\end{equation}
The curve is vertical when $\gamma$ is constant.  It is
$\boldsymbol\omega_i$-horizontal when
$v_{\boldsymbol\omega_i}(\Gamma)=0$.  It is mixed when $\dot\gamma$ and
$v_{\boldsymbol\omega_i}(\Gamma)$ both occur: equivalently, $\gamma$ is
nonconstant and $v_{\boldsymbol\omega_i}(\Gamma)$ is not identically zero.
Verticality is intrinsic to the bundle projection;
horizontality and the mixed decomposition depend on the selected connection.
\status{DEFINITION}

The definition agrees with the standard connection splitting of a bundle
tangent space \citep{Kobayashi1963}.  Since
$T\varpi_i\circ H^{\boldsymbol\omega_i}$ is the identity and
$T\varpi_i\circ\operatorname{ver}^{\boldsymbol\omega_i}=0$, applying
$T\varpi_i$ to \eqref{eq:hist-total-curve-splitting} recovers
$\dot\gamma$.  Thus a vertical curve stays in one fiber, whereas a
horizontal lift is a connection-selected comparison of points in different
fibers.  Neither operation moves the fixed base itself.
\status{ESTABLISHED}

\propositionheading{Horizontality is connection relative}{prop:hist-horizontal-connection-dependence}
The same nonvertical total-space curve can be mixed for one connection and
horizontal for another.  \status{ESTABLISHED}

\paragraph{Proof.}
Take the trivial line bundle
$\R_x\times\R_y\to\R_x$ and the curve
$\Gamma(r)=(r,r)$.  For the horizontal distribution
$H_0=\operatorname{span}\{\partial_x\}$, its vertical velocity is
$\partial_y$, so the curve is mixed.  For
$H_1=\operatorname{span}\{\partial_x+\partial_y\}$, the same velocity is
horizontal.  The base trace $x=r$ is unchanged.  $\square$

\section{Curves of sections and pointwise vertical change}
\label{sec:hist-section-curves}

\hypothesisheading{Regular section configuration}{hyp:hist-regular-section-space}
Let $\mathfrak S_i$ be a declared regular space of pairs of smooth sections
$Q=(q,s)$ of $\mathcal E_i\to\mathcal C_i$.  Whenever a tangent-space
calculation is used, assume either a finite-dimensional configuration
submanifold or a specified smooth locally convex manifold structure for
which evaluation is differentiable.  No such structure follows merely from
writing down all smooth sections.  \status{HYPOTHESIS}

For a $C^1$ representative $Q_i:I\to\mathfrak S_i$, define its adjoint
family by
\begin{equation}
\widehat\Sigma_i:I\times\mathcal C_i\longrightarrow\mathcal E_i,
\qquad
\widehat\Sigma_i(\lambda,c)=Q_i(\lambda)(c),
\qquad
\varpi_i\circ\widehat\Sigma_i(\lambda,c)=c.
\label{eq:hist-adjoint-section-family}
\end{equation}
Differentiation in the history parameter at fixed $c$ gives
\begin{equation}
T\varpi_i\left(\partial_\lambda\widehat\Sigma_i(\lambda,c)\right)=0,
\qquad
\partial_\lambda\widehat\Sigma_i(\lambda,c)
\in V_{Q_i(\lambda)(c)}\mathcal E_i.
\label{eq:hist-pointwise-history-verticality}
\end{equation}
Thus a curve of full sections produces a vertical statistical update at
every fixed context, even though its state is an entire section.
\status{ESTABLISHED}

At each history position the instantaneous sections also induce a pair of
connection-relative perceived geometries on the same fixed base:
\begin{equation}
\mathbf h_i(\lambda)
=\left(
h_{q_i(\lambda)}^{\omega_b},
h_{s_i(\lambda)}^{\omega_m}
\right).
\label{eq:hist-perceived-geometry-family}
\end{equation}
Consequently a section history determines an ordered family of base
semimetrics without making the base itself evolve.  Under
$\widetilde Q_i=Q_i\circ\phi$, this family changes only by
$\widetilde{\mathbf h}_i=\mathbf h_i\circ\phi$.  The connections are held
fixed in this statement; allowing them to change would add a separately
declared connection history and the correction terms of
\Cref{prop:pb-pullback-connection-change}. \status{ESTABLISHED}

\definitionheading{Oriented unparameterized history}{def:hist-oriented-history}
Two regular representatives $Q_i:I\to\mathfrak S_i$ and
$\widetilde Q_i:\widetilde I\to\mathfrak S_i$ represent the same oriented
history when
\begin{equation}
\widetilde Q_i=Q_i\circ\phi,
\qquad
\phi:\widetilde I\longrightarrow I
\text{ is an orientation-preserving $C^1$ diffeomorphism}.
\label{eq:hist-oriented-reparameterization-equivalence}
\end{equation}
The resulting abstract oriented curve is denoted by $\mathscr H_i$.  Its
adjoint evaluation has the required type
\begin{equation}
\Sigma_i:\mathscr H_i\times\mathcal C_i
\longrightarrow
\mathcal E_b\times_{\mathcal C}\mathcal E_m,
\qquad
\varpi_i\circ\Sigma_i(r,c)=c.
\label{eq:hist-adjoint-unparameterized-history}
\end{equation}
If a representative is an embedding, $\mathscr H_i$ may be identified with
its oriented image.  With self-intersections, $\mathscr H_i$ retains
traversal order and multiplicity and is not merely the set-theoretic image.
\status{DEFINITION}

The distinction between a base probe and a section history is exposed by a
diagonal evaluation.  Let $\gamma:I\to\mathcal C_i$ and set
$\Gamma(\lambda)=Q_i(\lambda)(\gamma(\lambda))$.  Write
\begin{equation}
D^{\boldsymbol\omega_i}Q_i(\lambda)
=\left(D^{\omega_b}q_i(\lambda),D^{\omega_m}s_i(\lambda)\right).
\label{eq:hist-product-covariant-section-jet}
\end{equation}
The chain rule and the connection splitting give
\begin{equation}
v_{\boldsymbol\omega_i}(\Gamma)
=\partial_\lambda\widehat\Sigma_i(\lambda,\gamma(\lambda))
+D^{\boldsymbol\omega_i}Q_i(\lambda)\dot\gamma(\lambda).
\label{eq:hist-diagonal-evaluation-velocity}
\end{equation}
The first term is history change at a fixed context.  The second is the
vertical response of one instantaneous section to a changing contextual
probe.  If $\gamma$ is constant, only the pointwise history velocity remains;
if $Q_i$ is constant, only the section jet remains.  When $\gamma$ is
constant, the diagonal curve is vertical regardless of whether the remaining
vertical velocity vanishes.  On an open subinterval where $\dot\gamma\ne0$,
the diagonal curve is mixed exactly when the sum in
\eqref{eq:hist-diagonal-evaluation-velocity} is nonzero, and it is horizontal
when that sum vanishes. \status{ESTABLISHED}

A smooth section is injective because
$\varpi_i\circ Q=\operatorname{id}_{\mathcal C_i}$, and it is an embedding
under the usual smooth-bundle hypotheses.  A total-space curve can therefore
be written as $Q\circ\gamma$ exactly when it lies in the section image, in
which case $\gamma=\varpi_i\circ\Gamma$ is unique.  For a general statistical
parameterization that is not a section, such descent is only local under a
local embedding hypothesis and is globally unique only under global
injectivity.  In particular, a nonconstant curve confined to one statistical
fiber cannot be the pushforward of a nonconstant base curve through a
section.  \status{ESTABLISHED}

\section{VFE selection of an oriented orbit}
\label{sec:hist-vfe-orbit-selection}

\hypothesisheading{Regular metric inference domain}{hyp:hist-regular-metric-domain}
Let $\mathcal Q_i\subseteq\mathfrak S_i$ be a regular configuration manifold
equipped with a positive-definite Fisher Riemannian metric
$\mathsf G_i^F$.  In an infinite-dimensional realization, assume it is a
strong metric so that the musical map is invertible on the covectors used
below.  Let $\Fenergy_i:\mathcal Q_i\to\R$ be $C^2$.  A quotient-metric
domain may replace $\mathcal Q_i$ only after the quotient has been
constructed.  \status{HYPOTHESIS}

\hypothesisheading{Exact multi-agent VFE lift}{hyp:hist-exact-vfe-lift}
Calling the configuration objective an exact local or collective VFE requires
more than its pair of marginal sections.  Fix a block $B$, an outside marginal
$Q_{B^c}$, and a declared manifold $\mathfrak R_B$ of admissible conditional
recognition kernels.  There must be a declared smooth marginalization or
configuration-extraction map
\begin{equation}
\pi_i^{\mathrm{conf}}:\mathfrak R_B\longrightarrow\mathcal Q_i
\label{eq:hist-configuration-extraction}
\end{equation}
on the admissible family.  Supply a smooth right-inverse lift
\begin{equation}
\iota_i:\mathcal Q_i\longrightarrow\mathfrak R_B,
\qquad
Q_i\longmapsto r_B^{Q_i}(\mathord{\cdot}\given b),
\qquad
\pi_i^{\mathrm{conf}}\circ\iota_i=\operatorname{id}_{\mathcal Q_i},
\label{eq:hist-exact-vfe-lift}
\end{equation}
so the lifted conditional law actually reproduces the displayed agent
configuration.  Equivalently, the full recognition law
$Q_{B^c}(db)r_B^{Q_i}(dy_B\given b)$ must have the declared marginals or
section statistics encoded by $Q_i$.  Without this commuting condition,
$\iota_i$ is merely a parameterized family of laws, not a lift of the agent
configuration.

For every $Q_i$ in the inference domain, impose the support and finiteness
hypotheses of \Cref{thm:obs-local-multiagent-elbo,thm:obs-local-global-potential}:
$Q_{B^c}\ll\Pi_{o,B^c}$,
$Z_B\in(0,+\infty)$ $Q_{B^c}$-almost everywhere,
$r_B^{Q_i}(\cdot\given b)\ll P_{0,B}(\cdot\given b)$, and integrability of the
KL terms and positive and negative energy parts, so the displayed local,
collective, and difference functionals are finite.  Assume also that the
lifted family is $C^2$ and dominated sufficiently to differentiate twice under
the $Q_{B^c}$ expectation.  Then define the outside-averaged objective
\begin{equation}
\overline{\Fenergy}_{B,o}(Q_i)
=\E_{Q_{B^c}}
\left[\Fenergy_{B,o}
(r_B^{Q_i};Y_{B^c})\right].
\label{eq:hist-outside-averaged-local-vfe}
\end{equation}
The lift is responsible for all correlation or copula data not determined by
the displayed marginal sections.  A fixed product or fixed-copula family is a
special hypothesis that can supply it; the pair $(q,s)$ alone cannot.
Alternatively, a population configuration may be lifted directly to a full
recognition law and the collective VFE pulled back along that lift.
If $\mathsf G_i^F$ is called the exact recognition Fisher metric, also require
\begin{equation}
\mathsf G_i^F=\iota_i^*G_{\mathfrak R_B}^F
\label{eq:hist-exact-fisher-lift}
\end{equation}
on a nondegenerate tier, or use its justified metric quotient.  A weighted
product of marginal-fiber metrics equals this joint pullback only under a
separately proved block-orthogonality or fixed-dependence hypothesis.
\status{HYPOTHESIS}

By \Cref{thm:obs-local-global-potential},
$\overline{\Fenergy}_{B,o}$ differs from the collective VFE restricted along
$Q_i\mapsto Q_{B^c}r_B^{Q_i}$ only by a $Q_i$-independent outside term.
Their differentials and natural-gradient rays therefore agree when the same
block metric is used.  A single realized conditional objective
$\Fenergy_{B,o}(r_B;b)$ need not have that gradient, and independently moving
outside marginals do not define the same fixed objective. \status{ESTABLISHED}

The Fisher natural gradient is characterized by
\begin{equation}
\mathsf G_i^F\left(\operatorname{grad}^F\Fenergy_i,Z\right)
=d\Fenergy_i[Z]
\quad\text{for every }Z\in T\mathcal Q_i.
\label{eq:hist-fisher-gradient-definition}
\end{equation}
At a noncritical configuration $Q$, define the positive descent ray
\begin{equation}
\mathcal R^-_{\Fenergy_i}(Q)
=\left\{-a\operatorname{grad}^F\Fenergy_i(Q):a>0\right\}.
\label{eq:hist-vfe-descent-ray}
\end{equation}
A natural-gradient history for the declared objective is an oriented
unparameterized history whose regular representatives satisfy
\begin{equation}
\dot Q_i(\lambda)
=-a_i(\lambda)\operatorname{grad}^F\Fenergy_i(Q_i(\lambda)),
\qquad a_i(\lambda)>0.
\label{eq:hist-vfe-ray-representative}
\end{equation}
The objective alone does not select this ray: the metric, or an explicitly
declared mobility replacing it, is load-bearing structure.
When $\Fenergy_i$ is the pullback specified in
\Cref{hyp:hist-exact-vfe-lift}, this is an exact VFE history; otherwise it is a
history for a generic declared objective and no exact-VFE identification is
asserted. \status{DEFINITION}

\propositionheading{Scalar mobility preserves the oriented orbit}{prop:hist-scalar-mobility-orbit}
On a noncritical domain, multiplication of the natural-gradient vector field
by any smooth positive scalar mobility changes only its parameterization.
A positive-definite anisotropic mobility can change the unparameterized
orbit.  \status{ESTABLISHED}

\paragraph{Proof.}
If $X=-\operatorname{grad}^F\Fenergy_i$ and $X_a=aX$ with $a>0$, then the
integral curves have the same tangent line and orientation.  Along an
$X$-integral curve, the change of parameter defined by a positive integral of
$a^{-1}$ converts one differential equation into the other.  For the second
claim, use the Euclidean plane with
\begin{equation}
\Fenergy(x,y)=\frac{x^2+y^2}{2},
\qquad
M=\begin{pmatrix}1&0\\0&2\end{pmatrix}.
\label{eq:hist-anisotropic-mobility-example}
\end{equation}
Ordinary gradient descent from $(x_0,y_0)$ satisfies
$y=(y_0/x_0)x$, whereas
$(\dot x,\dot y)^\top=-M\nabla\Fenergy$ satisfies
\begin{equation}
y=y_0\left(\frac{x}{x_0}\right)^2.
\label{eq:hist-anisotropic-orbit-change}
\end{equation}
For generic nonzero initial data these are different images.  $\square$

A VFE natural-gradient history is not generally a Fisher geodesic.  On the
same Euclidean Fisher plane, choose
\begin{equation}
\Fenergy(x,y)=\frac{x^2+2y^2}{2}.
\label{eq:hist-natural-gradient-nongeodesic}
\end{equation}
Its natural-gradient flow obeys $\dot x=-x$ and $\dot y=-2y$, hence follows
the curved trajectory in \eqref{eq:hist-anisotropic-orbit-change}.  Euclidean
geodesics are straight lines, so gradient selection and geodesic selection
are distinct even in this elementary Fisher geometry.  \status{ESTABLISHED}

\propositionheading{A semidefinite tensor cannot be silently inverted}{prop:hist-semidefinite-gradient-obstruction}
If the proposed Fisher tensor on configuration space is semidefinite, the
natural-gradient equation can fail to have a solution or can have nonunique
solutions.  \status{ESTABLISHED}

\paragraph{Proof.}
On $\R^2$ let $\mathsf G=dx^2$.  For $\Fenergy=y$, the equation
$\mathsf G(\operatorname{grad}\Fenergy,\cdot)=dy$ has no solution because
the left side annihilates $\partial_y$.  For $\Fenergy=x$, every vector
$\partial_x+b\partial_y$ solves the equation.  A justified metric quotient
or a separately declared regularization is therefore required before an
inverse or pseudoinverse is used.  $\square$

\section{Fisher duration on a selected history}
\label{sec:hist-fisher-duration}

\definitionheading{Fisher length and information clock}{def:hist-fisher-clock}
For a $C^1$ representative $Q:I\to\mathcal Q_i$, define its Fisher speed,
length, and cumulative clock from a chosen origin $\lambda_0$ by
\begin{align}
\nu_F(\lambda)
&=\sqrt{\mathsf G_i^F(\dot Q(\lambda),\dot Q(\lambda))},
\label{eq:hist-fisher-speed}\\
L_F(Q;[a,b])
&=\int_a^b\nu_F(\lambda)d\lambda,
\label{eq:hist-fisher-length}\\
\tau_{Q,\lambda_0}(\lambda)
&=\int_{\lambda_0}^{\lambda}\nu_F(u)du.
\label{eq:hist-cumulative-fisher-clock}
\end{align}
The origin is relational data attached to the history.  The auxiliary
parameter $\lambda$ is not part of the resulting arc length.
\status{DEFINITION}

\theoremheading{Reparameterization invariance and unit-speed form}{thm:hist-fisher-clock-invariance}
Let $\widetilde Q=Q\circ\phi$, where
$\phi:\widetilde I\to I$ is an orientation-preserving $C^1$
diffeomorphism.  Then corresponding subarcs have the same Fisher length and
their cumulative clocks obey
\begin{equation}
\tau_{\widetilde Q,\widetilde\lambda_0}(\widetilde\lambda)
=\tau_{Q,\lambda_0}(\phi(\widetilde\lambda)),
\qquad
\lambda_0=\phi(\widetilde\lambda_0).
\label{eq:hist-clock-reparameterization-law}
\end{equation}
If $\nu_F>0$, the cumulative clock is strictly increasing and gives the
unit-speed representative of the oriented history.  If $Q$ satisfies
\eqref{eq:hist-vfe-ray-representative}, then on its noncritical part
\begin{align}
\frac{d\tau}{d\lambda}
&=a_i\left\|\operatorname{grad}^F\Fenergy_i\right\|_F,
\label{eq:hist-vfe-clock-rate}\\
\frac{dQ}{d\tau}
&=-\frac{\operatorname{grad}^F\Fenergy_i}
{\left\|\operatorname{grad}^F\Fenergy_i\right\|_F},
\label{eq:hist-vfe-unit-tangent}\\
\frac{d\Fenergy_i}{d\tau}
&=-\left\|\operatorname{grad}^F\Fenergy_i\right\|_F.
\label{eq:hist-vfe-dissipation-per-length}
\end{align}
\status{ESTABLISHED}

\paragraph{Proof.}
The chain rule gives
$\dot{\widetilde Q}=(\dot Q\circ\phi)\phi'$ and hence
$\widetilde\nu_F=(\nu_F\circ\phi)\phi'$ because $\phi'>0$.
The substitution $u=\phi(\widetilde\lambda)$ proves both length invariance
and \eqref{eq:hist-clock-reparameterization-law}.  Positive speed makes
$\tau$ locally and globally invertible on the interval, and division by
$d\tau/d\lambda$ gives unit speed.  Substitution of
\eqref{eq:hist-vfe-ray-representative} into
\eqref{eq:hist-fisher-speed} gives
\eqref{eq:hist-vfe-clock-rate}.  The remaining identities follow from
\eqref{eq:hist-fisher-gradient-definition}.  $\square$

This is the standard Riemannian arc-length construction applied to the
selected statistical configuration manifold
\citep{AmariNagaoka2000,Pennec2006}.  Fisher geometry supplies duration;
the VFE descent ray supplies orientation.  A null segment in a semimetric or
a critical stationary segment has $\nu_F=0$, so the cumulative clock stalls
and cannot be used as a coordinate there.  \status{ESTABLISHED}

\propositionheading{Endpoints do not determine an update duration}{prop:hist-realized-length-versus-distance}
Let $A$ and $B$ lie in one path component of a regular Fisher manifold.
Those endpoints alone determine neither an update curve nor its realized
duration.  For every actual update $Q$ from $A$ to $B$,
\begin{equation}
d_F(A,B)
=\inf_{\substack{\Gamma(a)=A\\\Gamma(b)=B}}L_F(\Gamma;[a,b])
\leq L_F(Q;[a,b]),
\label{eq:hist-distance-versus-realized-length}
\end{equation}
and the inequality can be strict.  \status{ESTABLISHED}

\paragraph{Proof.}
The first equality is the definition of Riemannian distance as the infimum
of curve lengths.  Since the realized update $Q$ is one admissible curve,
the inequality follows.  In the Euclidean Fisher geometry of a two-parameter
normal location family, a curved path between two points is longer than the
straight geodesic joining them, which proves possible strictness.  $\square$

The information clock is also distinct from the amount of auxiliary
algorithmic parameter used to approach a limit.  In the unit-variance normal
location family, the path $\mu(t)=e^{-t}$ has Fisher length
\begin{equation}
\int_0^\infty |\dot\mu(t)|dt
=\int_0^\infty e^{-t}dt
=1,
\label{eq:hist-finite-length-infinite-parameter}
\end{equation}
although $t$ ranges over an unbounded interval.  No physical-time
interpretation follows from the length construction.  \status{NOT-CLAIMED}

\section{Pointwise, design, continuum, and quotient clocks}
\label{sec:hist-clock-tiers}

At one fixed $c\in\mathcal C_i$, choose positive belief and model weights
$w_b(c)$ and $w_m(c)$.  The pointwise squared speed of a section history is
\begin{equation}
\nu_{i,c}^2
=w_b(c)g_b^F(\partial_\lambda q_i(c),\partial_\lambda q_i(c))
+w_m(c)g_m^F(\partial_\lambda s_i(c),\partial_\lambda s_i(c)).
\label{eq:hist-pointwise-clock-speed}
\end{equation}
Integration of $\nu_{i,c}$ defines a pointwise clock.  This uses only
vertical evaluation velocities from
\eqref{eq:hist-pointwise-history-verticality}; it does not integrate over
the base.  \status{DEFINITION}

For a finite design $D=\{c_a\}_{a=1}^M$, declare positive design weights
$\rho_a$ and define
\begin{equation}
\nu_{i,D}^2
=\sum_{a=1}^M\rho_a\left[
w_b(c_a)g_b^F(\partial_\lambda q_i(c_a),\partial_\lambda q_i(c_a))
+w_m(c_a)g_m^F(\partial_\lambda s_i(c_a),\partial_\lambda s_i(c_a))
\right].
\label{eq:hist-finite-design-clock-speed}
\end{equation}
This is the product Fisher speed for the declared design and weights.
It is not automatically a positive-definite metric on the full section
space: any nonzero section variation vanishing at all $c_a$ lies in its
radical.  On a smooth positive-dimensional base, a bump variation supported
away from $D$ is an explicit witness whenever the vertical fiber has a
nonzero local direction there.  \status{ESTABLISHED}

A continuum clock requires a finite positive base measure $\mu_i$, measurable
positive channel weights, square-integrable evaluation velocities, and a
declared treatment of gauge redundancy.  Under those hypotheses its proposed
squared speed is
\begin{equation}
\nu_{i,\mu}^2
=\int_{\mathcal C_i}\left[
w_b(c)g_b^F(\partial_\lambda q_i(c),\partial_\lambda q_i(c))
+w_m(c)g_m^F(\partial_\lambda s_i(c),\partial_\lambda s_i(c))
\right]d\mu_i(c).
\label{eq:hist-continuum-clock-speed}
\end{equation}
The contextual base, the common principal bundle, and the VFE functional do
not by themselves select $\mu_i$, the relative channel weights, or a strong
manifold topology.  Without those declarations,
\eqref{eq:hist-continuum-clock-speed} is not an available clock and only the
pointwise or finite-design constructions are defined.  \status{ESTABLISHED}

On the finite-dimensional tier, let a finite-dimensional Lie group
$\mathcal G_i$ act freely, properly, and isometrically on the regular
Riemannian configuration manifold.  The quotient then carries the Riemannian
submersion metric.  Its speed is
\begin{equation}
\|[\dot Q]\|_{\mathrm{quot}}^2
=\inf_{\zeta\in T_Q(\mathcal G_i\mathbin{\cdot}Q)}
\mathsf G_i^F(\dot Q-\zeta,\dot Q-\zeta),
\label{eq:hist-quotient-gauge-speed}
\end{equation}
which equals the norm of the component orthogonal to the gauge orbit.  A
history-parameter-dependent coordinate representative can acquire spurious velocity
tangent to that orbit, while
\eqref{eq:hist-quotient-gauge-speed} is representative independent.  Passive
re-trivialization is a change of coordinates, not an additional physical
segment of the history.  If the action is not free and proper, the orbit
space can be stratified; a smooth quotient clock is then available only on a
declared regular stratum or after a separate singular construction.
\status{ESTABLISHED}

In an infinite-dimensional realization, free, proper, and isometric action
is not by itself enough for \eqref{eq:hist-quotient-gauge-speed}.  Require a
smooth principal quotient, closed split orbit-tangent subbundles, and smooth
orthogonal complements, equivalently a declared Riemannian-submersion or
mechanical-connection structure.  Otherwise the infimum can vanish on a
vector in the closure of a nonclosed orbit tangent without being attained,
so no orthogonal gauge quotient clock has been constructed.
\status{HYPOTHESIS}

\section{The obstruction to a global unit-speed clock}
\label{sec:hist-global-clock}

On a connected noncritical open set $U\subseteq\mathcal Q_i$, set
\begin{equation}
N=\left\|\operatorname{grad}^F\Fenergy_i\right\|_F,
\qquad
U_F=-\frac{\operatorname{grad}^F\Fenergy_i}{N},
\qquad
\alpha_F=U_F^\flat=-\frac{d\Fenergy_i}{N}.
\label{eq:hist-normalized-vfe-one-form}
\end{equation}
The form $\alpha_F$ evaluates to one on the unit descent field $U_F$ and
annihilates the Fisher-orthogonal complement of that field.
\status{DEFINITION}

\theoremheading{Exact criterion for a global orthogonal clock}{thm:hist-global-clock-exactness}
There is a smooth scalar $T:U\to\R$ satisfying
\begin{equation}
dT=\alpha_F,
\label{eq:hist-global-clock-potential}
\end{equation}
if and only if $\alpha_F$ is exact.  Equivalently, it must be closed and have
vanishing period around every piecewise $C^1$ closed curve in $U$:
\begin{equation}
d\alpha_F=0,
\qquad
\oint_\gamma\alpha_F=0
\quad\text{for every closed }\gamma\subset U.
\label{eq:hist-global-clock-period-criterion}
\end{equation}
On a simply connected $U$, closedness suffices.  Whenever $T$ exists, every
unit-speed VFE history in $U$ obeys $dT/d\tau=1$, and the level sets of $T$
are Fisher-orthogonal to the descent field.  \status{ESTABLISHED}

\paragraph{Proof.}
If \eqref{eq:hist-global-clock-potential} holds, exact forms are closed and
have zero periods.  Conversely, fix $Q_0\in U$ and define $T(Q)$ as the
integral of $\alpha_F$ from $Q_0$ to $Q$.  Closedness and vanishing periods
make this integral path independent, and the standard line-integral argument
gives $dT=\alpha_F$.  Since
$\alpha_F(U_F)=\mathsf G_i^F(U_F,U_F)=1$, the asserted unit rate follows.
The kernel of $dT=U_F^\flat$ is the orthogonal complement of $U_F$.
$\square$

The local differential obstruction is explicit:
\begin{equation}
d\alpha_F
=\frac{1}{N^2}dN\wedge d\Fenergy_i.
\label{eq:hist-normalized-form-curvature}
\end{equation}
Thus the norm of the gradient must be locally constant along each connected
regular level set of $\Fenergy_i$ for the normalized form to be closed.
Even after local closedness is established, the period condition in
\eqref{eq:hist-global-clock-period-criterion} is the remaining global
obligation.  \status{ESTABLISHED}

For a concrete local obstruction, take the Euclidean plane with
$\Fenergy(x,y)=xy$ on the first quadrant.  Then
$N=(x^2+y^2)^{1/2}$ and
\begin{equation}
d\alpha_F
=\frac{x^2-y^2}{(x^2+y^2)^{3/2}}dx\wedge dy,
\label{eq:hist-nonexact-clock-example}
\end{equation}
which is not identically zero on any open neighborhood.  Every individual
regular orbit still has its cumulative arc-length clock, but those orbitwise
origins do not assemble into the orthogonal scalar
\eqref{eq:hist-global-clock-potential}.  At critical points $N=0$ and
$\alpha_F$ is undefined, so extension across critical strata is a separate
problem.  The objective $\Fenergy_i$ itself may remain a Lyapunov scalar, but
it is not generally a unit Fisher clock.  \status{ESTABLISHED}

\section{Operational record clocks and exact Markov loss}
\label{sec:hist-record-clocks}

Let $\{P_\theta\}$ be a regular fine statistical experiment and let
$K:\mathsf X\rightsquigarrow\mathsf Y$ be a normalized,
parameter-independent Markov kernel.  Along a $C^1$ fine history
$\theta=\theta(\lambda)$, write $\ell^X_\lambda$ for its directional score
and define the record history by
\begin{equation}
P^Y_\lambda=P^X_\lambda K.
\label{eq:hist-same-path-markov-image}
\end{equation}
This equation identifies the record family as the image of the same fine
path; it does not define a separately optimized record trajectory.
\status{HYPOTHESIS}

\theoremheading{Same-path record and meta-agent clock contraction}{thm:hist-record-clock-contraction}
Under the preceding regularity and Markov hypotheses, the record score and
the exact Fisher-speed defect are
\begin{align}
\ell^Y_\lambda(Y)
&=\E_\lambda[\ell^X_\lambda(X)\mid Y],
\label{eq:hist-record-score-projection}\\
(\nu^X_F)^2-(\nu^Y_F)^2
&=\E_\lambda\operatorname{Var}_\lambda
(\ell^X_\lambda(X)\mid Y)
\geq0.
\label{eq:hist-record-conditional-variance-defect}
\end{align}
Consequently, on every common parameter interval,
\begin{equation}
L_F(P^Y)\leq L_F(P^X),
\qquad
\tau_Y(\lambda_1)-\tau_Y(\lambda_0)
\leq
\tau_X(\lambda_1)-\tau_X(\lambda_0).
\label{eq:hist-record-clock-contraction}
\end{equation}
Equality holds exactly when the fine directional score is measurable with
respect to the record for almost every path position.  The same result
applies to a meta-agent only when its statistical law is the normalized
parameter-independent Markov image of this same fine history.
\status{ESTABLISHED}

\paragraph{Proof.}
Attach $K$ to the fine experiment.  Because $K$ has no $\theta$ dependence,
the joint score is $\ell^X_\lambda(X)$.  Projection onto the record sigma
algebra gives \eqref{eq:hist-record-score-projection}.  The law of total
variance gives \eqref{eq:hist-record-conditional-variance-defect}, including
its equality condition.  Taking nonnegative square roots and integrating
proves \eqref{eq:hist-record-clock-contraction}.  This is the pathwise form
of the Fisher score-projection theorem
\citep{AyJostLeSchwachhoefer2018}.  $\square$

Parameter independence is load bearing.  If a dominated channel has density
$k_\lambda(y\mid x)$, its output score is instead
\begin{equation}
\ell^Y_\lambda(Y)
=\E_\lambda\left[
\ell^X_\lambda(X)
+\partial_\lambda\log k_\lambda(Y\mid X)
\mathrel{\big|}Y
\right].
\label{eq:hist-parameter-dependent-channel-score}
\end{equation}
The additional term can create Fisher information.  For example, a
parameter-independent input combined with a channel that emits a
Bernoulli variable of success probability
$(1+e^{-\lambda})^{-1}$ has zero fine Fisher speed and positive output Fisher
speed.  A fitted coarse approximation, an energy precomposition, a Galerkin
restriction, or a mismatch between generative and recognition parameters is
likewise outside
\Cref{thm:hist-record-clock-contraction}.  \status{ESTABLISHED}

\begin{figure}[htbp]
\centering
\begin{tikzpicture}[>=stealth]
  \node[draw,rounded corners,fill=lightblue,minimum width=3.3cm,minimum height=0.9cm]
    (fine) at (0,3.2) {fine history $P_r$};
  \node[draw,rounded corners,fill=lightgreen,minimum width=3.8cm,minimum height=0.9cm]
    (image) at (6.2,3.2) {same-path image $P_rK$};
  \draw[->,very thick,primaryblue] (fine)--(image);
  \node[align=center,text=primaryblue,font=\small,text width=4.8cm]
    at (3.1,4.15) {parameter-independent\\Markov channel};
  \node[align=center,font=\small,text width=6.4cm] at (3.1,2.05)
    {exact conditional-variance defect; $d\tau_{P K}\leq d\tau_P$};
  \node[draw,rounded corners,fill=lightblue,minimum width=3.3cm,minimum height=0.9cm]
    (fflow) at (0,0) {$\Phi_t(Q_0)$};
  \node[draw,rounded corners,fill=lightgreen,minimum width=3.8cm,minimum height=0.9cm]
    (mflow) at (6.2,0) {$\bar\Phi_u(\mathcal RQ_0)$};
  \draw[->,very thick,dashed,sciencegreen] (fflow)--(mflow);
  \node[align=center,text=sciencegreen,font=\small,text width=4.8cm]
    at (3.1,0.95) {independent flows require\\oriented semiconjugacy};
  \node[align=center,font=\small,text width=6.3cm] at (3.1,-1.05)
    {$\mathcal C$ is fixed; the RG level $\ell$ is depth, not the history parameter $r$};
\end{tikzpicture}
\caption{The upper arrow is the same-path comparison covered by Fisher
contraction.  The lower arrow compares independently generated flows and
exists only after the additional semiconjugacy condition is proved.}
\label{fig:hist-markov-clock-and-semiconjugacy}
\end{figure}

\section{Independently recomputed flows}
\label{sec:hist-independent-flows}

Let $X_f$ and $X_m$ be the declared fine and meta-agent natural-gradient
vector fields on regular configuration manifolds $\mathcal Q_f$ and
$\mathcal Q_m$, with local flows $\Phi_t$ and $\bar\Phi_u$.  Let
$\mathcal R:\mathcal Q_f\to\mathcal Q_m$ be a smooth coarse map.  Pointwise
contraction of $T\mathcal R$ does not state that $\mathcal R$ maps an
$X_f$-orbit to an $X_m$-orbit.  \status{ESTABLISHED}

\propositionheading{Oriented semiconjugacy criterion}{prop:hist-oriented-semiconjugacy}
On a noncritical domain, $\mathcal R$ maps each fine integral curve to a meta
integral curve up to an orientation-preserving change of parameter if
\begin{equation}
T_Q\mathcal R X_f(Q)
=a(Q)X_m(\mathcal RQ),
\qquad a(Q)>0.
\label{eq:hist-oriented-vector-field-semiconjugacy}
\end{equation}
More precisely, uniqueness of the two local flows gives
\begin{equation}
\mathcal R\circ\Phi_t(Q)
=\bar\Phi_{\sigma_Q(t)}(\mathcal RQ),
\qquad
\sigma_Q(t)=\int_0^t a(\Phi_s(Q))ds.
\label{eq:hist-oriented-flow-semiconjugacy}
\end{equation}
Conversely, differentiating an oriented flow semiconjugacy of this form at
$t=0$ yields \eqref{eq:hist-oriented-vector-field-semiconjugacy}.
\status{ESTABLISHED}

\paragraph{Proof.}
The derivative of $\mathcal R\circ\Phi_t(Q)$ is
$a(\Phi_tQ)X_m(\mathcal R\Phi_tQ)$ by
\eqref{eq:hist-oriented-vector-field-semiconjugacy}.  Since
$\dot\sigma_Q=a(\Phi_tQ)>0$, the right side of
\eqref{eq:hist-oriented-flow-semiconjugacy} solves the same initial-value
problem after the stated time change.  Flow uniqueness proves the identity.
The converse is the chain rule at $t=0$.  $\square$

For a meta-agent clock comparison between independently recomputed flows,
two distinct facts are therefore required.  The coarse map must be a
parameter-independent Markov pushforward to invoke
\Cref{thm:hist-record-clock-contraction}, and it must satisfy
\eqref{eq:hist-oriented-vector-field-semiconjugacy} to identify the
independently recomputed meta orbit with the image of the fine orbit.  Either
condition without the other is insufficient.  At a fine critical point,
the same relation with $a>0$ also requires the image to be meta-critical.
Establishing these conditions for the manuscript's independently recomputed
RG flows remains an open obligation.  \status{OPEN}

\section{RG depth is not inference time}
\label{sec:hist-rg-depth-not-time}

A hierarchy of coarse objects may be indexed by
\begin{equation}
Q^{(\ell+1)}=\mathcal R_\ell(Q^{(\ell)}),
\qquad \ell=0,1,2,\ldots.
\label{eq:hist-rg-depth-index}
\end{equation}
The integer $\ell$ records RG depth.  It is not a parameter on an inference
history, and the sequence in \eqref{eq:hist-rg-depth-index} is not a temporal
trajectory.  A history can instead be present at every level,
$Q^{(\ell)}(r)$, with $r$ labeling positions on its oriented orbit.  Only
when
\begin{equation}
Q^{(\ell+1)}(r)
=\mathcal R_\ell(Q^{(\ell)}(r)),
\label{eq:hist-levelwise-same-path-image}
\end{equation}
and $\mathcal R_\ell$ is a parameter-independent Markov map does the
same-path clock contraction compare adjacent levels.  Independently
recomputed levelwise histories additionally require
\Cref{prop:hist-oriented-semiconjugacy}.  \status{ESTABLISHED}

Cross-scale covariant first jets, their connection-mismatch term, and the
additive Fisher-defect cocycle belong to
\Cref{ch:pullback-geometry}.  They compare typed geometry across a scale map;
they do not turn the scale index into time.  The construction in this chapter
therefore yields oriented, parameter-free inference histories and
agent-relative information duration while leaving any identification with
physical time outside the established theory.  \status{NOT-CLAIMED}
